\chapter{Introdução}

A incorporação de ferramentas de geração de código assistida por inteligência artificial (IA) ao cotidiano do desenvolvimento de software ocorreu de forma acelerada e abrangente. Assistentes como GitHub Copilot, ChatGPT e similares passaram, em poucos anos, de curiosidades tecnológicas a componentes efetivos do fluxo de trabalho de milhões de desenvolvedores ao redor do mundo. O argumento central que impulsiona essa adoção é a produtividade: experimentos controlados demonstram que programadores assistidos por IA concluem tarefas de codificação até 55\% mais rapidamente do que seus pares sem assistência \cite{peng2023}. Ferramentas de suporte à decisão automatizadas têm como objetivo declarado, em diversas áreas, ``melhorar a confiabilidade e a qualidade da tomada de decisão humana no trabalho ou na vida cotidiana'' \cite{mosier1998} — e as ferramentas de geração de código se enquadram precisamente nessa categoria.

Esse ganho, no entanto, não ocorre de forma neutra. À medida que a geração automática de código se torna mais fluente e convincente, emerge uma questão que este trabalho se propõe a investigar: a rapidez com que o código é produzido corresponde à qualidade e à correção do que é entregue? A resposta da literatura recente é perturbadora. \citeonline{pearce2022}, ao avaliar as contribuições de código do GitHub Copilot em 89 cenários baseados nas vulnerabilidades mais críticas catalogadas pelo MITRE CWE Top 25, constataram que aproximadamente 40\% do código gerado continha falhas de segurança. Mais revelador ainda é o estudo de \citeonline{perry2023}: em experimento com desenvolvedores reais realizando tarefas de segurança, o grupo assistido por IA produziu código menos seguro em quatro de cinco tarefas — e o fez com maior confiança em suas próprias soluções do que o grupo sem assistência.

Essa combinação — mais erros, menos percepção dos erros — não é acidental. Ela encontra explicação em um mecanismo cognitivo amplamente documentado na psicologia social, anterior ao surgimento da IA generativa em décadas: o conceito de \textit{cognitive miser}, introduzido por \citeonline{fiske1984}. Segundo esse modelo, o ser humano é, por padrão, um agente avaro em termos de esforço cognitivo: o cérebro recorre a heurísticas e atalhos mentais sempre que possível, reservando o raciocínio analítico profundo apenas para situações em que não há alternativa. Como sintetiza \citeonline{mosier1998}, ``a perspectiva do \textit{cognitive miser} sugere que as pessoas tendem a tomar o caminho de menor esforço cognitivo e podem difundir a responsabilidade por tarefas decisórias a auxílios automatizados, resultando em verificação cruzada incompleta das informações disponíveis''. A IA generativa não criou essa tendência — ela simplesmente oferece, pela primeira vez, uma via de menor esforço suficientemente boa para tarefas que antes exigiam obrigatoriamente atenção crítica, como escrever e revisar código.

O fenômeno resultante é conhecido como \textit{viés de automação}: a tendência de aceitar, sem questionamento, as sugestões de sistemas automatizados, mesmo quando o indivíduo possui conhecimento técnico suficiente para identificar problemas \cite{mosier1998}. O viés se manifesta em dois tipos de erro: erros de omissão, quando o operador deixa de agir porque o sistema não alertou; e erros de comissão, quando o operador age incorretamente por seguir uma recomendação equivocada do sistema. A gravidade desse mecanismo fica evidente em estudos com pilotos profissionais: diante de um alerta falso de incêndio gerado por um sistema automatizado, 75\% dos pilotos seguiram a recomendação errada do auxílio à decisão; em uma variante do experimento, 100\% desligaram o motor em resposta ao alerta errôneo — a ação mais severa e irreversível possível \cite{mosier1998}. Mais perturbador ainda: 67\% dos pilotos relataram ter visto, no painel, indicações adicionais que confirmavam o alerta de incêndio — indicações que simplesmente não existiam. Esse fenômeno, denominado \textit{phantom memory}, revela que o cérebro não apenas segue o sistema cegamente, mas constrói retroativamente memórias falsas para justificar a decisão tomada \cite{mosier1998}. Em contraste, apenas 3\% dos eventos críticos foram ignorados pelo grupo de controle que realizou a mesma tarefa sem suporte automatizado.

Importa notar que o \textit{viés de automação} não é o único padrão disfuncional na relação entre humanos e sistemas automatizados. Seu inverso, a \textit{aversão a algoritmos}, foi documentado por \citeonline{dietvorst2015}: pessoas abandonam algoritmos superiores após vê-los cometer um único erro, tolerando erros humanos equivalentes sem a mesma reação. Conforme demonstrado em cinco estudos experimentais, isso ocorre porque ``as pessoas perdem confiança em previsores algorítmicos mais rapidamente do que em previsores humanos ao vê-los cometer o mesmo erro'' — mesmo quando os algoritmos objetivamente superam seus pares humanos. Juntos, o viés de automação e a aversão a algoritmos delimitam o espectro psicológico dentro do qual o desenvolvedor contemporâneo opera ao utilizar ferramentas de IA: ora confia demais, ora desconfia erraticamente.

As consequências cognitivas de longo prazo do uso de ferramentas de IA são igualmente preocupantes. \citeonline{gerlich2025}, em estudo empírico com 666 participantes, encontrou correlação negativa significativa entre o uso frequente de ferramentas de IA e a capacidade de pensamento crítico, mediada pelo aumento do \textit{cognitive offloading} — a externalização de processos cognitivos para ferramentas externas a fim de reduzir a carga sobre a memória de trabalho. O estudo identificou o risco de ``dependência cognitiva, em que indivíduos tornam-se excessivamente dependentes de ferramentas de IA para tarefas rotineiras e complexas'', e observou que ``o problema da caixa-preta pode reduzir o engajamento crítico e a responsabilidade, à medida que indivíduos confiam cegamente nas recomendações da IA sem questioná-las ou avaliá-las''. Em contextos profissionais, esse padrão resulta em trabalhadores que ``podem ter dificuldade em tomar decisões de forma independente quando esses sistemas estão indisponíveis'' \cite{gerlich2025}.

Evidência experimental direta foi fornecida por \citeonline{georgiou2025}, em único experimento randomizado controlado medindo engajamento cognitivo: o grupo que utilizou ChatGPT apresentou escores significativamente menores de engajamento cognitivo em comparação ao grupo de controle. O achado mais alarmante diz respeito à persistência do efeito: participantes que migraram da IA para o trabalho sem assistência continuaram apresentando desempenho inferior mesmo após a remoção da ferramenta, evidenciando que a dependência cognitiva induzida pela IA deixa rastros mensuráveis mesmo quando o auxílio não está mais disponível. Complementarmente, \citeonline{lee2025} demonstraram, em estudo com trabalhadores do conhecimento, que ``maior confiança na IA generativa está associada a menor pensamento crítico, enquanto maior autoconfiança está associada a maior pensamento crítico'' — sugerindo que a deferência ao sistema substitui, e não complementa, o julgamento independente.

Agrava esse cenário o fato de que os próprios benchmarks utilizados para avaliar a qualidade dos modelos de geração de código tendem a superestimar seu desempenho. \citeonline{liu2023} demonstraram que o HumanEval — referência padrão da área — superestima a taxa de correção dos modelos em até 29\%, em razão de conjuntos de testes insuficientes. Na prática de engenharia de software real, a distância entre o que os modelos aparentam fazer e o que efetivamente entregam é ainda mais pronunciada: no \textit{benchmark} SWE-bench, que testa modelos em \textit{issues} reais de repositórios do GitHub, os maiores modelos disponíveis resolveram menos de 2\% dos casos \cite{jimenez2024}.

\section{Objetivos}

Este trabalho tem como objetivo geral analisar como o viés de automação e os mecanismos de descarga cognitiva influenciam a qualidade e a correção do código produzido com o auxílio de ferramentas de IA generativa, sob a perspectiva da Engenharia de Computação.

Como objetivos específicos, busca-se:

\begin{itemize}
	\item revisar os principais estudos empíricos que avaliam a qualidade, a segurança e a correção do código gerado por modelos de linguagem de grande porte;
	\item examinar os fundamentos teóricos da psicologia cognitiva — em especial o \textit{cognitive miser}, o viés de automação e a aversão a algoritmos — que explicam a relação disfuncional entre desenvolvedores e ferramentas de IA;
	\item avaliar as consequências cognitivas de longo prazo do uso de ferramentas de IA generativa, incluindo declínio de pensamento crítico e dependência cognitiva;
	\item identificar práticas e mecanismos de desenvolvimento de software que possam atuar como mitigadores do viés de automação no contexto da IA generativa;
	\item discutir as implicações desses achados para a formação e a prática profissional em Engenharia de Computação.
\end{itemize}

\section{Organização do Trabalho}

O presente trabalho está organizado da seguinte forma. O Capítulo 2 apresenta a revisão da literatura, abrangendo os fundamentos teóricos da cognição humana relevantes para o problema, os principais estudos empíricos sobre qualidade e segurança do código gerado por IA, e os mecanismos de avaliação utilizados na área. O Capítulo 3 descreve a metodologia adotada para a condução desta pesquisa. O Capítulo 4 apresenta e discute os resultados obtidos. Por fim, o Capítulo 5 expõe as conclusões do trabalho e aponta direções para pesquisas futuras.
